% =====================================================================
%  RELATED WORK
%  MSc thesis, Beirut Arab University
%  Ontology-based skincare recommendation for the Lebanese market
%
%  Include from your main file with:   % =====================================================================
%  RELATED WORK
%  MSc thesis, Beirut Arab University
%  Ontology-based skincare recommendation for the Lebanese market
%
%  Include from your main file with:   % =====================================================================
%  RELATED WORK
%  MSc thesis, Beirut Arab University
%  Ontology-based skincare recommendation for the Lebanese market
%
%  Include from your main file with:   % =====================================================================
%  RELATED WORK
%  MSc thesis, Beirut Arab University
%  Ontology-based skincare recommendation for the Lebanese market
%
%  Include from your main file with:   \input{related_work}
%  Requires in the preamble:
%     \usepackage{booktabs}
%     \usepackage{graphicx}
%     \usepackage{tikz}
%     \usetikzlibrary{positioning,arrows.meta,shapes.geometric}
%     \usepackage[hidelinks]{hyperref}
%     \usepackage{multirow}
%     \usepackage{rotating}      % for the wide comparison table
%     \usepackage{threeparttable}
%
%  Bibliography: related_work.bib
% =====================================================================

\chapter{Related Work}
\label{ch:related}

% ---------------------------------------------------------------------
\section{How this chapter is organised}
\label{sec:rw-organisation}

This chapter narrows in four steps. Section~\ref{sec:rw-datasets} reviews how
product datasets and semantic resources are constructed and what the community
expects of them, which positions the first contribution of this thesis.
Section~\ref{sec:rw-ontologies} reviews ontologies and knowledge graphs built
for skincare, cosmetics and cosmetic regulation, which are the direct
competitors. Section~\ref{sec:rw-kgrec} reviews ontology and knowledge graph
based recommender systems in other domains, from which the mechanisms of this
work are drawn. Section~\ref{sec:rw-dl} reviews deep learning approaches,
including those applied to skin analysis, and states plainly what they do
better. Section~\ref{sec:rw-hybrid} reviews hybrid neuro-symbolic work and
positions it as future rather than present work. Section~\ref{sec:rw-gap}
consolidates the gap.

Each section closes by naming what is missing, and the missing thing is what
the following section addresses. The final gap is the thesis.

\begin{figure}[htbp]
\centering
\begin{tikzpicture}[
  node distance=7mm,
  box/.style={draw, rounded corners, align=center, text width=0.82\linewidth,
              inner sep=5pt, font=\small},
  gap/.style={draw, thick, rounded corners, align=center,
              text width=0.82\linewidth, inner sep=5pt, font=\small\bfseries},
  arr/.style={-{Latex[length=2mm]}, thick}
]
\node[box] (a) {\textbf{\S\ref{sec:rw-datasets} Resource construction}\\
  provenance, validation and availability are expected of a modern resource};
\node[box, below=of a] (b) {\textbf{\S\ref{sec:rw-ontologies} Skincare and cosmetics ontologies}\\
  small, mostly unpublished, ingredients as free text, no regulatory grounding};
\node[box, below=of b] (c) {\textbf{\S\ref{sec:rw-kgrec} KG recommenders elsewhere}\\
  mature mechanisms exist, but none has been applied to skincare};
\node[box, below=of c] (d) {\textbf{\S\ref{sec:rw-dl} Deep learning}\\
  effective where labelled data exists, cannot cite a regulator, needs interactions};
\node[box, below=of d] (e) {\textbf{\S\ref{sec:rw-hybrid} Hybrid neuro-symbolic}\\
  assumes a domain knowledge graph already exists};
\node[gap,  below=of e] (f) {\S\ref{sec:rw-gap} Gap: no skincare knowledge graph with
  regulator-linked ingredients, modelled provenance and real market availability};
\draw[arr] (a) -- (b); \draw[arr] (b) -- (c);
\draw[arr] (c) -- (d); \draw[arr] (d) -- (e); \draw[arr] (e) -- (f);
\end{tikzpicture}
\caption{The structure of this chapter. Each section identifies a limitation
that the next addresses.}
\label{fig:rw-structure}
\end{figure}

% ---------------------------------------------------------------------
\section{Dataset and resource construction}
\label{sec:rw-datasets}

% TODO Lynne: paste your existing dataset related work from main.tex here,
% then keep the closing paragraph below, which is what links this section
% to the next one.

Resource contributions are an established category with their own venues and
their own review criteria. FoodKG was published as a resource paper at the
International Semantic Web Conference~\cite{haussmann2019foodkg}, and the TOXIN
knowledge graph appeared in \emph{Database}~\cite{toxin2024}. Both are judged on
availability, documentation, validation and a stated maintenance plan rather
than on predictive accuracy.

% TODO Lynne: one paragraph here on existing cosmetic product datasets and
% why none covers a national market. Use the material already in main.tex.

\paragraph{What is missing.} Existing cosmetic product resources record what a
product is. None records whether the product can be obtained in a given market,
what it costs there, or how much the source of each claim is worth. The datasets
that do exist are not linked to any regulatory register, so no statement they
support can be traced to a legal instrument. The next section shows that the
ontologies built over such data inherit all three limitations.

% ---------------------------------------------------------------------
\section{Ontologies and knowledge graphs for skincare and cosmetics}
\label{sec:rw-ontologies}

\subsection{Cross-domain cosmetics ontologies}
\label{ssec:rw-moe}

Moe and Aung~\cite{moe2014} construct two ontologies, one describing facial
skin problems and one describing cosmetics, joined by bridge relations. Their
methodological contribution is an explicit eight-task ontological engineering
procedure, running from a glossary of terms with synonyms and acronyms through
concept taxonomies, binary relation diagrams and a concept dictionary to a
constants table. Both ontologies were built in Prot\'eg\'e.

Recommendation proceeds in three stages. Taxonomic conversational case-based
reasoning narrows an initially vague user query through ranked questions, using
an \texttt{isNextRelatedTo} property to order them. The resulting problem and
the candidate products are then placed in a weighted directed acyclic graph, and
the Ford--Fulkerson maximum flow algorithm assigns each product a flow weight
$W(v_i)=\sum_k f(v_i,k)$, with products ranked in descending order of that
weight. Evaluation reports precision, recall and $F$-measure as a function of a
recommendation threshold $\alpha$.

Two observations matter for this thesis. First, no reasoner is used at any
point, so the ontology serves as a graph structure rather than as a logical
theory. Second, the paper claims greater accuracy than related work without
reporting a dataset size, a baseline or a comparison, so that claim is not
supported by the evidence presented. Ingredients are represented as a free-text
property with an associated numeric value and are not linked to any register,
which means no statement about regulatory status is expressible in the model.

\subsection{OntoCosmetic and formulation support}
\label{ssec:rw-ontocosmetic}

Serna et al.~\cite{serna2021} present OntoCosmetic, a decision support ontology
for the design of emulsion-based cosmetic products, extended by Gabriel et
al.~\cite{gabriel2023} into a cross-platform application, Formultools. The
knowledge base is assembled from four kinds of material: general subproblems
expressed as physicochemical properties to promote or limit, general solution
strategies, ingredient databases typed by function, and heuristics linking the
two. Heuristics are encoded as SWRL rules in Prot\'eg\'e and carry
\texttt{hasHeuristicHighThreshold}, \texttt{hasHeuristicLowThreshold} and
\texttt{hasHeuristicSource} properties, the last of which records the provenance
of the rule itself.

The application was developed using the five-planes user-centred design
method, in co-design with cosmetics experts, and evaluated across iterations
using the AttrakDiff instrument. It supports ingredient screening, multi-criteria
ingredient selection and formulation evaluation, citing the Analytic Hierarchy
Process for the multi-criteria component.

OntoCosmetic is the only ontology reviewed in this section that is publicly
resolvable, at \url{https://purl.org/ontocosmetic}. Inspection of the published
file gives 116 classes, 26 object properties, 20 data properties and 279
individuals, all populated manually. Its classes, among them \texttt{HLB},
\texttt{DropletSize}, \texttt{Rheology} and \texttt{Emollient\_Dosage}, model
formulation chemistry rather than retail selection: the ontology is designed to
help a formulator make a stable cream, not to help a consumer choose one.
Neither paper reports an evaluation against ground truth, and Gabriel et al.
state in their conclusion that testing with experts remained future work.

\subsection{Vision-driven skincare recommendation}
\label{ssec:rw-hansanie}

Hansanie and Silva~\cite{hansanie2024} combine a convolutional neural network
that grades acne severity from a facial photograph with an ontology that
performs the product selection. Their vocabulary was derived from interviews
with dermatologists and health professionals together with a survey of
twenty-one participants, and the class hierarchy was built top-down. The
ontology is deliberately split into three modules, covering skincare concepts,
product information and user profiles, which are then merged; the stated
rationale is that the three evolve independently.

Their property set includes \texttt{suitableFor} from treatment products to
skin types, \texttt{hasKeyIngredient}, \texttt{hasProductRecommendation} and
\texttt{hasRating}. The implementation uses Prot\'eg\'e, the Pellet reasoner,
Owlready2 for querying from Python, and Tkinter for the interface. The severity
grade produced by the network is asserted into the ontology as a fact about the
user, so the perception component and the knowledge component remain separable.

The reported figures require care. The headline result of 87.5\,\% is the
proportion of twenty-four survey participants who expressed satisfaction with
the products shown; the classification accuracy of the network itself is
77.5\,\%. No baseline recommender is compared against, and no ground-truth set
of correct recommendations is constructed.

\subsection{Defined classes and query-driven recommendation}
\label{ssec:rw-abesova}

Abesov\'a et al.~\cite{abesova2023} build a skincare ontology over Sephora
product and review data, comprising 36 classes, 6 object properties and 6 data
properties, following an explicitly middle-out and iterative construction
process in which the scope widened twice as new data sources were incorporated.
The knowledge graph was produced from scraped CSV files using OntoRefine and
stored in GraphDB, with country data and capital-city coordinates obtained by
querying DBpedia through an external SPARQL endpoint. Recommendation is
performed by SPARQL queries rather than by procedural code.

Their principal methodological contribution, and the one adopted in this
thesis, is the use of defined classes. Rather than annotating individual
products with a suitability label, they state an equivalence axiom, so that
membership of \texttt{OilySkinProducts} follows from
\texttt{hasOilyScore\ value\ 5} and is recomputed by the reasoner whenever the
underlying data changes.

The work is a student project and is not peer reviewed, which is stated here
because it is cited for a mechanism rather than for a result. Its three
acknowledged limitations are directly relevant. Defined-class membership
provides no ordering within a class, so ranking requires a separate signal; a
product supported by a single weak review can be recommended ahead of products
with strong support, because the two selection paths are not commensurable; and
the authors observe that class restrictions based on ingredients would permit a
symbolic rather than statistical approach, which they did not implement.

\subsection{Ontologies populated from commercial catalogues}
\label{ssec:rw-bittech}

A recent ontology-based skincare recommendation system~\cite{bittech2025}
follows METHONTOLOGY to produce twelve core classes, among them Product,
Ingredient, Skin Type, Skin Concern, Benefit, Allergen Type and Formulation
Trait, together with more than twenty-five object properties. The knowledge
base was populated by scraping three commercial platforms, yielding over 3{,}800
products and 28{,}000 ingredients, and is queried through Apache Jena Fuseki.

This is the closest published system to the present work in both scale and
intent. It differs in four respects. It targets the Indonesian market rather
than a market in which product availability is itself uncertain; its allergen
representation is a locally defined class rather than a link to a statutory
register; it does not model price, availability or provenance; and its product
count is approximately one third of the dataset used here.

% TODO Lynne: obtain and read the full PDF before submission, then confirm or
% correct the four points above and add whether their ontology is published.

\subsection{Regulatory and clinical knowledge graphs}
\label{ssec:rw-regulatory}

Two lines of work outside recommendation are relevant because they supply the
grounding that skincare ontologies lack.

The TOXIN knowledge graph~\cite{toxin2024} supports animal-free risk assessment
of cosmetics, drawing on the scientific opinions that underpin the annexes of
Regulation (EC) No~1223/2009. Its construction is instructive: curated
spreadsheets are transformed into RDF using R2RML, the ToXic Process Ontology
from the OBO Foundry provides the primary structure, external resources are
incorporated by reference to their IRIs rather than by copying, and each
integration is held in a separate named graph so that the origin of any
statement can be traced. It contains no products.

In the clinical direction, DermO~\cite{dermo2016} provides more than three
thousand expert-authored terms for cutaneous disease, organised under twenty
upper-level categories, aligned to ICD-10 and published through BioPortal in
both OBO and OWL~2 formats. The more recent D3X ontology~\cite{d3x2024} links
dermoscopic patterns to differential diagnoses and is evaluated through a
usability study. Neither addresses cosmetic products, but both offer a
vocabulary for skin concerns with clinical standing that a retail-derived
vocabulary does not possess.

A parallel may also be drawn with ontologies that model regulatory status
directly. Work on flavouring ontologies for halal status determination
\cite{halal2024} represents the ingredient, the certifying authority and the
resulting status as distinct entities, so that a product inherits a status from
its constituents according to a named authority. The structure transfers
directly to ingredients restricted under the European cosmetics regulation.

\subsection{Summary of Section~\ref{sec:rw-ontologies}}

\begin{table}[htbp]
\centering
\caption{Ontologies and knowledge graphs in skincare, cosmetics and cosmetic
regulation. A dash indicates the property is absent or not reported.}
\label{tab:rw-ontologies}
\footnotesize
\setlength{\tabcolsep}{4pt}
\begin{tabular}{lccccccc}
\toprule
 & \rotatebox{60}{Moe 2014} & \rotatebox{60}{OntoCosmetic}
 & \rotatebox{60}{Hansanie 2024} & \rotatebox{60}{Abesov\'a 2023}
 & \rotatebox{60}{Bit-Tech 2025} & \rotatebox{60}{TOXIN 2024}
 & \rotatebox{60}{\textbf{This work}} \\
\midrule
Peer reviewed                & $\circ$ & \checkmark & \checkmark & --         & \checkmark & \checkmark & \checkmark \\
Ontology published           & --      & \checkmark & --         & --         & --         & \checkmark & \checkmark \\
Named methodology            & $\circ$ & --         & --         & \checkmark & \checkmark & --         & \checkmark \\
Products modelled            & --      & 279        & $\circ$    & $\circ$    & 3{,}800    & --         & \textbf{12{,}629} \\
Ingredients linked to registry & --    & --         & --         & --         & --         & \checkmark & \checkmark \\
Declarative population       & --      & --         & --         & \checkmark & --         & \checkmark & \checkmark \\
Reasoner used                & --      & \checkmark & \checkmark & \checkmark & \checkmark & --         & \checkmark \\
SPARQL interface             & --      & --         & $\circ$    & \checkmark & \checkmark & \checkmark & \checkmark \\
External linking             & --      & --         & --         & \checkmark & --         & \checkmark & \checkmark \\
Provenance modelled          & --      & $\circ$    & --         & --         & --         & \checkmark & \checkmark \\
Price or availability        & $\circ$ & $\circ$    & --         & $\circ$    & --         & --         & \checkmark \\
Evaluated on ground truth    & --      & --         & $\circ$    & --         & $\circ$    & n/a        & \checkmark \\
\bottomrule
\end{tabular}
\begin{tablenotes}\footnotesize
\item \checkmark\ present; $\circ$ partial or weak; -- absent or not reported.
\end{tablenotes}
\end{table}

\paragraph{What is missing.} No system in Table~\ref{tab:rw-ontologies}
combines an ingredient layer with statutory grounding, an explicit model of
where each claim came from, and a representation of whether the recommended
product can actually be obtained. Ingredients are represented as free text in
every skincare ontology reviewed; the single work that links ingredients to the
European regulatory framework contains no products at all. The mechanisms
needed to close this gap have, however, been developed in other domains, and
these are reviewed next.

% ---------------------------------------------------------------------
\section{Ontology and knowledge graph based recommender systems}
\label{sec:rw-kgrec}

\subsection{Ontological user profiling}

Middleton et al.~\cite{middleton2004} established ontological user profiling in
their Quickstep and Foxtrot systems for recommending research papers, reporting
three findings that remain the standard justification for the approach:
ontological inference improves user profiling, because interest in a specific
topic implies interest in its parent; external ontological knowledge can
bootstrap a recommender for users with no history; and visualising a profile so
that users may correct it improves accuracy. The second finding is the direct
answer to the cold-start problem that affects any purely collaborative method,
and it applies with particular force to a market in which no interaction logs
exist.

\subsection{Constraint satisfaction over a domain knowledge graph}

FoodKG~\cite{haussmann2019foodkg} integrates recipes, nutritional data, food
taxonomies and links to existing ontologies into a single graph, over which a
SPARQL service identifies recipes that can be prepared from available
ingredients while respecting hard constraints such as allergies. The structural
correspondence with the present work is close: recipes correspond to products,
nutritional data supplied by an authority corresponds to the CosIng register,
and allergen constraints correspond to the twenty-six fragrance allergens whose
declaration is required under the European regulation. Subsequent work frames
the task explicitly as constrained question answering over a knowledge
graph~\cite{chen2021constrained}, which is a more accurate description of the
skincare selection problem than ranking is.

\subsection{Linked open data as a source of item features}

Di Noia, Ostuni and colleagues~\cite{dinoia2012lod,ostuni2013topn,dinoia2015wims}
compute item similarity using a semantic vector space model whose dimensions are
links in the linked open data cloud rather than terms in a document. The
approach permits content-based recommendation in the absence of ratings, which
is the situation of this thesis, where the rating field is populated for only
half of the catalogue and no user history exists at all. The authors identify
entity resolution against external datasets as their principal open problem, an
observation consistent with the matching error rates measured during the
construction of the dataset used here.

\subsection{User needs as first-class entities}

AliCoCo~\cite{alicoco2020} argues that conventional product ontologies leave a
semantic gap because they describe what a product is rather than what a user
needs, and formalises e-commerce concepts, that is user needs, as first-class
entities in their own right. The argument transfers directly. Skin concerns are
presently modelled throughout the literature as attributes of products, whereas
a need entity can be satisfied by a combination of products, can carry a budget
and a season, and can be constrained by local availability.

\subsection{Taxonomies of knowledge graph recommendation}

Guo et al.~\cite{guo2022survey} classify knowledge-graph-based recommender
systems as embedding-based, path-based or unified. All three families presuppose
a user--item interaction matrix and treat the knowledge graph as side
information mitigating sparsity. The present work is situated upstream of this
taxonomy: it constructs, for a domain and a market in which none exists, the
knowledge graph that such methods presuppose.

\paragraph{What is missing.} The mechanisms surveyed in this section, defined
classes, hard constraint satisfaction, linked-data-derived similarity and needs
as entities, are mature and well evidenced, but none has been applied to
skincare, and none of the systems reviewed operates in a market where
availability cannot be assumed.

% ---------------------------------------------------------------------
\section{Deep learning approaches}
\label{sec:rw-dl}

Lee et al.~\cite{lee2024} estimate the efficacy of a cosmetic product from its
ingredient list using a deep neural network, combining the result with automated
facial skin analysis to produce personalised recommendations. The point of
departure is the same as that of this thesis, namely the ingredient list, and
the divergence lies entirely in how ingredients are connected to outcomes: by a
function learned from data in their case, and by the European register together
with stated rules in this one.

The broader literature on neural recommendation is well
surveyed~\cite{zhang2019survey}, with neural collaborative
filtering~\cite{he2017ncf} providing the standard baseline. Every method in this
family learns from user behaviour. The dataset constructed for this thesis
contains none, because it is a product resource rather than an interaction log,
and this determines which half of the literature the work belongs to.

It should be stated directly that where sufficient labelled data exists, a
learned model is likely to outperform a rule-based one on benchmark accuracy,
because it can capture interactions that are not explicitly encoded. The
argument advanced here is not one of accuracy. It is that in a domain carrying
physical risk and governed by a legal framework, a recommendation that cannot
state its reason cannot be audited by a regulator, cannot be used by a
pharmacist, and cannot be corrected when it is wrong.

\paragraph{What is missing.} Vision and learned models can establish the state
of a user's skin and can estimate what a formulation is likely to do. Neither
can establish that the resulting product is stocked in Beirut, what it costs
there this month, or which annex of the cosmetics regulation restricts one of
its preservatives.

% ---------------------------------------------------------------------
\section{Hybrid neuro-symbolic approaches}
\label{sec:rw-hybrid}

A substantial body of work combines knowledge graphs with neural
recommendation. CKE~\cite{zhang2016cke} learns item representations jointly from
a knowledge graph and from ratings; RippleNet~\cite{wang2018ripplenet} propagates
user preference outward through the graph; KGAT~\cite{wang2019kgat} applies graph
attention over a collaborative knowledge graph to capture high-order relations;
and KPRN~\cite{wang2019kprn} encodes the path between user and item with a
recurrent network and weights paths by their strength, so that the path itself
constitutes an explanation. Subsequent work optimises explanation paths for
quality rather than relevance alone~\cite{balloccu2022}.

Each of these methods requires an interaction matrix and treats the knowledge
graph as auxiliary. The position of this thesis is the converse: the graph is
the contribution and the interactions do not yet exist. These methods are
therefore what the resource constructed here enables once deployment produces
interaction data, and they are discussed again in the future work of
Chapter~\ref{ch:conclusion}.

% TODO Lynne: check that \ref{ch:conclusion} matches your actual label.

One bridge is applicable immediately. Ontology embedding methods such as
RDF2Vec~\cite{ristoski2016rdf2vec} and OWL2Vec*~\cite{chen2021owl2vec} require an
ontology rather than an interaction log. OWL2Vec* in particular encodes graph
structure, lexical content and logical constructors, having first materialised
the ontology with a reasoner, and was evaluated on class membership and class
subsumption prediction. Both tasks correspond to real gaps in the present
dataset, in which size is recorded for 27\,\% of products and rating for 51\,\%.

Finally, the use of large language models in ontology engineering is developing
rapidly, with a dedicated shared task~\cite{llms4ol2024} covering term typing,
taxonomy discovery and non-taxonomic relation extraction, and a recent
systematic review of the area~\cite{li2026llmoe}. The free-text benefit and
concern fields of the dataset used here are candidates for exactly these tasks,
subject to validation of every proposed term against an authoritative source.

% ---------------------------------------------------------------------
\section{The gap addressed by this thesis}
\label{sec:rw-gap}

Five limitations recur across the reviewed literature.

\begin{enumerate}
  \item \textbf{No ingredient layer with statutory standing.} Every skincare
    ontology reviewed represents ingredients as free text or as a small,
    manually curated class list. The one knowledge graph that links cosmetic
    ingredients to the European regulatory framework contains no products.
  \item \textbf{No model of provenance.} No reviewed skincare ontology records
    who asserted a claim, or how much weight that source should carry. Every
    claim is presented as equally reliable.
  \item \textbf{No model of availability.} Every reviewed system assumes that a
    recommended product can be purchased. In the Lebanese market that assumption
    fails, and it fails differently for imported and for locally manufactured
    products.
  \item \textbf{Weak evaluation.} Two of the reviewed systems report user
    satisfaction on samples of fewer than thirty participants, one presents a
    single case study, and one states that expert evaluation was never carried
    out.
  \item \textbf{Poor availability of artefacts.} One of six reviewed ontologies
    is publicly resolvable, which prevents the field from building on its own
    results.
\end{enumerate}

This thesis addresses limitations 1, 2, 3 and 5 directly through the design of
the ontology and the publication of the resulting artefacts, and addresses
limitation 4 through a competency-question-driven evaluation reported in
Chapter~\ref{ch:evaluation}.

% TODO Lynne: check that \ref{ch:evaluation} matches your actual label.

%  Requires in the preamble:
%     \usepackage{booktabs}
%     \usepackage{graphicx}
%     \usepackage{tikz}
%     \usetikzlibrary{positioning,arrows.meta,shapes.geometric}
%     \usepackage[hidelinks]{hyperref}
%     \usepackage{multirow}
%     \usepackage{rotating}      % for the wide comparison table
%     \usepackage{threeparttable}
%
%  Bibliography: related_work.bib
% =====================================================================

\chapter{Related Work}
\label{ch:related}

% ---------------------------------------------------------------------
\section{How this chapter is organised}
\label{sec:rw-organisation}

This chapter narrows in four steps. Section~\ref{sec:rw-datasets} reviews how
product datasets and semantic resources are constructed and what the community
expects of them, which positions the first contribution of this thesis.
Section~\ref{sec:rw-ontologies} reviews ontologies and knowledge graphs built
for skincare, cosmetics and cosmetic regulation, which are the direct
competitors. Section~\ref{sec:rw-kgrec} reviews ontology and knowledge graph
based recommender systems in other domains, from which the mechanisms of this
work are drawn. Section~\ref{sec:rw-dl} reviews deep learning approaches,
including those applied to skin analysis, and states plainly what they do
better. Section~\ref{sec:rw-hybrid} reviews hybrid neuro-symbolic work and
positions it as future rather than present work. Section~\ref{sec:rw-gap}
consolidates the gap.

Each section closes by naming what is missing, and the missing thing is what
the following section addresses. The final gap is the thesis.

\begin{figure}[htbp]
\centering
\begin{tikzpicture}[
  node distance=7mm,
  box/.style={draw, rounded corners, align=center, text width=0.82\linewidth,
              inner sep=5pt, font=\small},
  gap/.style={draw, thick, rounded corners, align=center,
              text width=0.82\linewidth, inner sep=5pt, font=\small\bfseries},
  arr/.style={-{Latex[length=2mm]}, thick}
]
\node[box] (a) {\textbf{\S\ref{sec:rw-datasets} Resource construction}\\
  provenance, validation and availability are expected of a modern resource};
\node[box, below=of a] (b) {\textbf{\S\ref{sec:rw-ontologies} Skincare and cosmetics ontologies}\\
  small, mostly unpublished, ingredients as free text, no regulatory grounding};
\node[box, below=of b] (c) {\textbf{\S\ref{sec:rw-kgrec} KG recommenders elsewhere}\\
  mature mechanisms exist, but none has been applied to skincare};
\node[box, below=of c] (d) {\textbf{\S\ref{sec:rw-dl} Deep learning}\\
  effective where labelled data exists, cannot cite a regulator, needs interactions};
\node[box, below=of d] (e) {\textbf{\S\ref{sec:rw-hybrid} Hybrid neuro-symbolic}\\
  assumes a domain knowledge graph already exists};
\node[gap,  below=of e] (f) {\S\ref{sec:rw-gap} Gap: no skincare knowledge graph with
  regulator-linked ingredients, modelled provenance and real market availability};
\draw[arr] (a) -- (b); \draw[arr] (b) -- (c);
\draw[arr] (c) -- (d); \draw[arr] (d) -- (e); \draw[arr] (e) -- (f);
\end{tikzpicture}
\caption{The structure of this chapter. Each section identifies a limitation
that the next addresses.}
\label{fig:rw-structure}
\end{figure}

% ---------------------------------------------------------------------
\section{Dataset and resource construction}
\label{sec:rw-datasets}

% TODO Lynne: paste your existing dataset related work from main.tex here,
% then keep the closing paragraph below, which is what links this section
% to the next one.

Resource contributions are an established category with their own venues and
their own review criteria. FoodKG was published as a resource paper at the
International Semantic Web Conference~\cite{haussmann2019foodkg}, and the TOXIN
knowledge graph appeared in \emph{Database}~\cite{toxin2024}. Both are judged on
availability, documentation, validation and a stated maintenance plan rather
than on predictive accuracy.

% TODO Lynne: one paragraph here on existing cosmetic product datasets and
% why none covers a national market. Use the material already in main.tex.

\paragraph{What is missing.} Existing cosmetic product resources record what a
product is. None records whether the product can be obtained in a given market,
what it costs there, or how much the source of each claim is worth. The datasets
that do exist are not linked to any regulatory register, so no statement they
support can be traced to a legal instrument. The next section shows that the
ontologies built over such data inherit all three limitations.

% ---------------------------------------------------------------------
\section{Ontologies and knowledge graphs for skincare and cosmetics}
\label{sec:rw-ontologies}

\subsection{Cross-domain cosmetics ontologies}
\label{ssec:rw-moe}

Moe and Aung~\cite{moe2014} construct two ontologies, one describing facial
skin problems and one describing cosmetics, joined by bridge relations. Their
methodological contribution is an explicit eight-task ontological engineering
procedure, running from a glossary of terms with synonyms and acronyms through
concept taxonomies, binary relation diagrams and a concept dictionary to a
constants table. Both ontologies were built in Prot\'eg\'e.

Recommendation proceeds in three stages. Taxonomic conversational case-based
reasoning narrows an initially vague user query through ranked questions, using
an \texttt{isNextRelatedTo} property to order them. The resulting problem and
the candidate products are then placed in a weighted directed acyclic graph, and
the Ford--Fulkerson maximum flow algorithm assigns each product a flow weight
$W(v_i)=\sum_k f(v_i,k)$, with products ranked in descending order of that
weight. Evaluation reports precision, recall and $F$-measure as a function of a
recommendation threshold $\alpha$.

Two observations matter for this thesis. First, no reasoner is used at any
point, so the ontology serves as a graph structure rather than as a logical
theory. Second, the paper claims greater accuracy than related work without
reporting a dataset size, a baseline or a comparison, so that claim is not
supported by the evidence presented. Ingredients are represented as a free-text
property with an associated numeric value and are not linked to any register,
which means no statement about regulatory status is expressible in the model.

\subsection{OntoCosmetic and formulation support}
\label{ssec:rw-ontocosmetic}

Serna et al.~\cite{serna2021} present OntoCosmetic, a decision support ontology
for the design of emulsion-based cosmetic products, extended by Gabriel et
al.~\cite{gabriel2023} into a cross-platform application, Formultools. The
knowledge base is assembled from four kinds of material: general subproblems
expressed as physicochemical properties to promote or limit, general solution
strategies, ingredient databases typed by function, and heuristics linking the
two. Heuristics are encoded as SWRL rules in Prot\'eg\'e and carry
\texttt{hasHeuristicHighThreshold}, \texttt{hasHeuristicLowThreshold} and
\texttt{hasHeuristicSource} properties, the last of which records the provenance
of the rule itself.

The application was developed using the five-planes user-centred design
method, in co-design with cosmetics experts, and evaluated across iterations
using the AttrakDiff instrument. It supports ingredient screening, multi-criteria
ingredient selection and formulation evaluation, citing the Analytic Hierarchy
Process for the multi-criteria component.

OntoCosmetic is the only ontology reviewed in this section that is publicly
resolvable, at \url{https://purl.org/ontocosmetic}. Inspection of the published
file gives 116 classes, 26 object properties, 20 data properties and 279
individuals, all populated manually. Its classes, among them \texttt{HLB},
\texttt{DropletSize}, \texttt{Rheology} and \texttt{Emollient\_Dosage}, model
formulation chemistry rather than retail selection: the ontology is designed to
help a formulator make a stable cream, not to help a consumer choose one.
Neither paper reports an evaluation against ground truth, and Gabriel et al.
state in their conclusion that testing with experts remained future work.

\subsection{Vision-driven skincare recommendation}
\label{ssec:rw-hansanie}

Hansanie and Silva~\cite{hansanie2024} combine a convolutional neural network
that grades acne severity from a facial photograph with an ontology that
performs the product selection. Their vocabulary was derived from interviews
with dermatologists and health professionals together with a survey of
twenty-one participants, and the class hierarchy was built top-down. The
ontology is deliberately split into three modules, covering skincare concepts,
product information and user profiles, which are then merged; the stated
rationale is that the three evolve independently.

Their property set includes \texttt{suitableFor} from treatment products to
skin types, \texttt{hasKeyIngredient}, \texttt{hasProductRecommendation} and
\texttt{hasRating}. The implementation uses Prot\'eg\'e, the Pellet reasoner,
Owlready2 for querying from Python, and Tkinter for the interface. The severity
grade produced by the network is asserted into the ontology as a fact about the
user, so the perception component and the knowledge component remain separable.

The reported figures require care. The headline result of 87.5\,\% is the
proportion of twenty-four survey participants who expressed satisfaction with
the products shown; the classification accuracy of the network itself is
77.5\,\%. No baseline recommender is compared against, and no ground-truth set
of correct recommendations is constructed.

\subsection{Defined classes and query-driven recommendation}
\label{ssec:rw-abesova}

Abesov\'a et al.~\cite{abesova2023} build a skincare ontology over Sephora
product and review data, comprising 36 classes, 6 object properties and 6 data
properties, following an explicitly middle-out and iterative construction
process in which the scope widened twice as new data sources were incorporated.
The knowledge graph was produced from scraped CSV files using OntoRefine and
stored in GraphDB, with country data and capital-city coordinates obtained by
querying DBpedia through an external SPARQL endpoint. Recommendation is
performed by SPARQL queries rather than by procedural code.

Their principal methodological contribution, and the one adopted in this
thesis, is the use of defined classes. Rather than annotating individual
products with a suitability label, they state an equivalence axiom, so that
membership of \texttt{OilySkinProducts} follows from
\texttt{hasOilyScore\ value\ 5} and is recomputed by the reasoner whenever the
underlying data changes.

The work is a student project and is not peer reviewed, which is stated here
because it is cited for a mechanism rather than for a result. Its three
acknowledged limitations are directly relevant. Defined-class membership
provides no ordering within a class, so ranking requires a separate signal; a
product supported by a single weak review can be recommended ahead of products
with strong support, because the two selection paths are not commensurable; and
the authors observe that class restrictions based on ingredients would permit a
symbolic rather than statistical approach, which they did not implement.

\subsection{Ontologies populated from commercial catalogues}
\label{ssec:rw-bittech}

A recent ontology-based skincare recommendation system~\cite{bittech2025}
follows METHONTOLOGY to produce twelve core classes, among them Product,
Ingredient, Skin Type, Skin Concern, Benefit, Allergen Type and Formulation
Trait, together with more than twenty-five object properties. The knowledge
base was populated by scraping three commercial platforms, yielding over 3{,}800
products and 28{,}000 ingredients, and is queried through Apache Jena Fuseki.

This is the closest published system to the present work in both scale and
intent. It differs in four respects. It targets the Indonesian market rather
than a market in which product availability is itself uncertain; its allergen
representation is a locally defined class rather than a link to a statutory
register; it does not model price, availability or provenance; and its product
count is approximately one third of the dataset used here.

% TODO Lynne: obtain and read the full PDF before submission, then confirm or
% correct the four points above and add whether their ontology is published.

\subsection{Regulatory and clinical knowledge graphs}
\label{ssec:rw-regulatory}

Two lines of work outside recommendation are relevant because they supply the
grounding that skincare ontologies lack.

The TOXIN knowledge graph~\cite{toxin2024} supports animal-free risk assessment
of cosmetics, drawing on the scientific opinions that underpin the annexes of
Regulation (EC) No~1223/2009. Its construction is instructive: curated
spreadsheets are transformed into RDF using R2RML, the ToXic Process Ontology
from the OBO Foundry provides the primary structure, external resources are
incorporated by reference to their IRIs rather than by copying, and each
integration is held in a separate named graph so that the origin of any
statement can be traced. It contains no products.

In the clinical direction, DermO~\cite{dermo2016} provides more than three
thousand expert-authored terms for cutaneous disease, organised under twenty
upper-level categories, aligned to ICD-10 and published through BioPortal in
both OBO and OWL~2 formats. The more recent D3X ontology~\cite{d3x2024} links
dermoscopic patterns to differential diagnoses and is evaluated through a
usability study. Neither addresses cosmetic products, but both offer a
vocabulary for skin concerns with clinical standing that a retail-derived
vocabulary does not possess.

A parallel may also be drawn with ontologies that model regulatory status
directly. Work on flavouring ontologies for halal status determination
\cite{halal2024} represents the ingredient, the certifying authority and the
resulting status as distinct entities, so that a product inherits a status from
its constituents according to a named authority. The structure transfers
directly to ingredients restricted under the European cosmetics regulation.

\subsection{Summary of Section~\ref{sec:rw-ontologies}}

\begin{table}[htbp]
\centering
\caption{Ontologies and knowledge graphs in skincare, cosmetics and cosmetic
regulation. A dash indicates the property is absent or not reported.}
\label{tab:rw-ontologies}
\footnotesize
\setlength{\tabcolsep}{4pt}
\begin{tabular}{lccccccc}
\toprule
 & \rotatebox{60}{Moe 2014} & \rotatebox{60}{OntoCosmetic}
 & \rotatebox{60}{Hansanie 2024} & \rotatebox{60}{Abesov\'a 2023}
 & \rotatebox{60}{Bit-Tech 2025} & \rotatebox{60}{TOXIN 2024}
 & \rotatebox{60}{\textbf{This work}} \\
\midrule
Peer reviewed                & $\circ$ & \checkmark & \checkmark & --         & \checkmark & \checkmark & \checkmark \\
Ontology published           & --      & \checkmark & --         & --         & --         & \checkmark & \checkmark \\
Named methodology            & $\circ$ & --         & --         & \checkmark & \checkmark & --         & \checkmark \\
Products modelled            & --      & 279        & $\circ$    & $\circ$    & 3{,}800    & --         & \textbf{12{,}629} \\
Ingredients linked to registry & --    & --         & --         & --         & --         & \checkmark & \checkmark \\
Declarative population       & --      & --         & --         & \checkmark & --         & \checkmark & \checkmark \\
Reasoner used                & --      & \checkmark & \checkmark & \checkmark & \checkmark & --         & \checkmark \\
SPARQL interface             & --      & --         & $\circ$    & \checkmark & \checkmark & \checkmark & \checkmark \\
External linking             & --      & --         & --         & \checkmark & --         & \checkmark & \checkmark \\
Provenance modelled          & --      & $\circ$    & --         & --         & --         & \checkmark & \checkmark \\
Price or availability        & $\circ$ & $\circ$    & --         & $\circ$    & --         & --         & \checkmark \\
Evaluated on ground truth    & --      & --         & $\circ$    & --         & $\circ$    & n/a        & \checkmark \\
\bottomrule
\end{tabular}
\begin{tablenotes}\footnotesize
\item \checkmark\ present; $\circ$ partial or weak; -- absent or not reported.
\end{tablenotes}
\end{table}

\paragraph{What is missing.} No system in Table~\ref{tab:rw-ontologies}
combines an ingredient layer with statutory grounding, an explicit model of
where each claim came from, and a representation of whether the recommended
product can actually be obtained. Ingredients are represented as free text in
every skincare ontology reviewed; the single work that links ingredients to the
European regulatory framework contains no products at all. The mechanisms
needed to close this gap have, however, been developed in other domains, and
these are reviewed next.

% ---------------------------------------------------------------------
\section{Ontology and knowledge graph based recommender systems}
\label{sec:rw-kgrec}

\subsection{Ontological user profiling}

Middleton et al.~\cite{middleton2004} established ontological user profiling in
their Quickstep and Foxtrot systems for recommending research papers, reporting
three findings that remain the standard justification for the approach:
ontological inference improves user profiling, because interest in a specific
topic implies interest in its parent; external ontological knowledge can
bootstrap a recommender for users with no history; and visualising a profile so
that users may correct it improves accuracy. The second finding is the direct
answer to the cold-start problem that affects any purely collaborative method,
and it applies with particular force to a market in which no interaction logs
exist.

\subsection{Constraint satisfaction over a domain knowledge graph}

FoodKG~\cite{haussmann2019foodkg} integrates recipes, nutritional data, food
taxonomies and links to existing ontologies into a single graph, over which a
SPARQL service identifies recipes that can be prepared from available
ingredients while respecting hard constraints such as allergies. The structural
correspondence with the present work is close: recipes correspond to products,
nutritional data supplied by an authority corresponds to the CosIng register,
and allergen constraints correspond to the twenty-six fragrance allergens whose
declaration is required under the European regulation. Subsequent work frames
the task explicitly as constrained question answering over a knowledge
graph~\cite{chen2021constrained}, which is a more accurate description of the
skincare selection problem than ranking is.

\subsection{Linked open data as a source of item features}

Di Noia, Ostuni and colleagues~\cite{dinoia2012lod,ostuni2013topn,dinoia2015wims}
compute item similarity using a semantic vector space model whose dimensions are
links in the linked open data cloud rather than terms in a document. The
approach permits content-based recommendation in the absence of ratings, which
is the situation of this thesis, where the rating field is populated for only
half of the catalogue and no user history exists at all. The authors identify
entity resolution against external datasets as their principal open problem, an
observation consistent with the matching error rates measured during the
construction of the dataset used here.

\subsection{User needs as first-class entities}

AliCoCo~\cite{alicoco2020} argues that conventional product ontologies leave a
semantic gap because they describe what a product is rather than what a user
needs, and formalises e-commerce concepts, that is user needs, as first-class
entities in their own right. The argument transfers directly. Skin concerns are
presently modelled throughout the literature as attributes of products, whereas
a need entity can be satisfied by a combination of products, can carry a budget
and a season, and can be constrained by local availability.

\subsection{Taxonomies of knowledge graph recommendation}

Guo et al.~\cite{guo2022survey} classify knowledge-graph-based recommender
systems as embedding-based, path-based or unified. All three families presuppose
a user--item interaction matrix and treat the knowledge graph as side
information mitigating sparsity. The present work is situated upstream of this
taxonomy: it constructs, for a domain and a market in which none exists, the
knowledge graph that such methods presuppose.

\paragraph{What is missing.} The mechanisms surveyed in this section, defined
classes, hard constraint satisfaction, linked-data-derived similarity and needs
as entities, are mature and well evidenced, but none has been applied to
skincare, and none of the systems reviewed operates in a market where
availability cannot be assumed.

% ---------------------------------------------------------------------
\section{Deep learning approaches}
\label{sec:rw-dl}

Lee et al.~\cite{lee2024} estimate the efficacy of a cosmetic product from its
ingredient list using a deep neural network, combining the result with automated
facial skin analysis to produce personalised recommendations. The point of
departure is the same as that of this thesis, namely the ingredient list, and
the divergence lies entirely in how ingredients are connected to outcomes: by a
function learned from data in their case, and by the European register together
with stated rules in this one.

The broader literature on neural recommendation is well
surveyed~\cite{zhang2019survey}, with neural collaborative
filtering~\cite{he2017ncf} providing the standard baseline. Every method in this
family learns from user behaviour. The dataset constructed for this thesis
contains none, because it is a product resource rather than an interaction log,
and this determines which half of the literature the work belongs to.

It should be stated directly that where sufficient labelled data exists, a
learned model is likely to outperform a rule-based one on benchmark accuracy,
because it can capture interactions that are not explicitly encoded. The
argument advanced here is not one of accuracy. It is that in a domain carrying
physical risk and governed by a legal framework, a recommendation that cannot
state its reason cannot be audited by a regulator, cannot be used by a
pharmacist, and cannot be corrected when it is wrong.

\paragraph{What is missing.} Vision and learned models can establish the state
of a user's skin and can estimate what a formulation is likely to do. Neither
can establish that the resulting product is stocked in Beirut, what it costs
there this month, or which annex of the cosmetics regulation restricts one of
its preservatives.

% ---------------------------------------------------------------------
\section{Hybrid neuro-symbolic approaches}
\label{sec:rw-hybrid}

A substantial body of work combines knowledge graphs with neural
recommendation. CKE~\cite{zhang2016cke} learns item representations jointly from
a knowledge graph and from ratings; RippleNet~\cite{wang2018ripplenet} propagates
user preference outward through the graph; KGAT~\cite{wang2019kgat} applies graph
attention over a collaborative knowledge graph to capture high-order relations;
and KPRN~\cite{wang2019kprn} encodes the path between user and item with a
recurrent network and weights paths by their strength, so that the path itself
constitutes an explanation. Subsequent work optimises explanation paths for
quality rather than relevance alone~\cite{balloccu2022}.

Each of these methods requires an interaction matrix and treats the knowledge
graph as auxiliary. The position of this thesis is the converse: the graph is
the contribution and the interactions do not yet exist. These methods are
therefore what the resource constructed here enables once deployment produces
interaction data, and they are discussed again in the future work of
Chapter~\ref{ch:conclusion}.

% TODO Lynne: check that \ref{ch:conclusion} matches your actual label.

One bridge is applicable immediately. Ontology embedding methods such as
RDF2Vec~\cite{ristoski2016rdf2vec} and OWL2Vec*~\cite{chen2021owl2vec} require an
ontology rather than an interaction log. OWL2Vec* in particular encodes graph
structure, lexical content and logical constructors, having first materialised
the ontology with a reasoner, and was evaluated on class membership and class
subsumption prediction. Both tasks correspond to real gaps in the present
dataset, in which size is recorded for 27\,\% of products and rating for 51\,\%.

Finally, the use of large language models in ontology engineering is developing
rapidly, with a dedicated shared task~\cite{llms4ol2024} covering term typing,
taxonomy discovery and non-taxonomic relation extraction, and a recent
systematic review of the area~\cite{li2026llmoe}. The free-text benefit and
concern fields of the dataset used here are candidates for exactly these tasks,
subject to validation of every proposed term against an authoritative source.

% ---------------------------------------------------------------------
\section{The gap addressed by this thesis}
\label{sec:rw-gap}

Five limitations recur across the reviewed literature.

\begin{enumerate}
  \item \textbf{No ingredient layer with statutory standing.} Every skincare
    ontology reviewed represents ingredients as free text or as a small,
    manually curated class list. The one knowledge graph that links cosmetic
    ingredients to the European regulatory framework contains no products.
  \item \textbf{No model of provenance.} No reviewed skincare ontology records
    who asserted a claim, or how much weight that source should carry. Every
    claim is presented as equally reliable.
  \item \textbf{No model of availability.} Every reviewed system assumes that a
    recommended product can be purchased. In the Lebanese market that assumption
    fails, and it fails differently for imported and for locally manufactured
    products.
  \item \textbf{Weak evaluation.} Two of the reviewed systems report user
    satisfaction on samples of fewer than thirty participants, one presents a
    single case study, and one states that expert evaluation was never carried
    out.
  \item \textbf{Poor availability of artefacts.} One of six reviewed ontologies
    is publicly resolvable, which prevents the field from building on its own
    results.
\end{enumerate}

This thesis addresses limitations 1, 2, 3 and 5 directly through the design of
the ontology and the publication of the resulting artefacts, and addresses
limitation 4 through a competency-question-driven evaluation reported in
Chapter~\ref{ch:evaluation}.

% TODO Lynne: check that \ref{ch:evaluation} matches your actual label.

%  Requires in the preamble:
%     \usepackage{booktabs}
%     \usepackage{graphicx}
%     \usepackage{tikz}
%     \usetikzlibrary{positioning,arrows.meta,shapes.geometric}
%     \usepackage[hidelinks]{hyperref}
%     \usepackage{multirow}
%     \usepackage{rotating}      % for the wide comparison table
%     \usepackage{threeparttable}
%
%  Bibliography: related_work.bib
% =====================================================================

\chapter{Related Work}
\label{ch:related}

% ---------------------------------------------------------------------
\section{How this chapter is organised}
\label{sec:rw-organisation}

This chapter narrows in four steps. Section~\ref{sec:rw-datasets} reviews how
product datasets and semantic resources are constructed and what the community
expects of them, which positions the first contribution of this thesis.
Section~\ref{sec:rw-ontologies} reviews ontologies and knowledge graphs built
for skincare, cosmetics and cosmetic regulation, which are the direct
competitors. Section~\ref{sec:rw-kgrec} reviews ontology and knowledge graph
based recommender systems in other domains, from which the mechanisms of this
work are drawn. Section~\ref{sec:rw-dl} reviews deep learning approaches,
including those applied to skin analysis, and states plainly what they do
better. Section~\ref{sec:rw-hybrid} reviews hybrid neuro-symbolic work and
positions it as future rather than present work. Section~\ref{sec:rw-gap}
consolidates the gap.

Each section closes by naming what is missing, and the missing thing is what
the following section addresses. The final gap is the thesis.

\begin{figure}[htbp]
\centering
\begin{tikzpicture}[
  node distance=7mm,
  box/.style={draw, rounded corners, align=center, text width=0.82\linewidth,
              inner sep=5pt, font=\small},
  gap/.style={draw, thick, rounded corners, align=center,
              text width=0.82\linewidth, inner sep=5pt, font=\small\bfseries},
  arr/.style={-{Latex[length=2mm]}, thick}
]
\node[box] (a) {\textbf{\S\ref{sec:rw-datasets} Resource construction}\\
  provenance, validation and availability are expected of a modern resource};
\node[box, below=of a] (b) {\textbf{\S\ref{sec:rw-ontologies} Skincare and cosmetics ontologies}\\
  small, mostly unpublished, ingredients as free text, no regulatory grounding};
\node[box, below=of b] (c) {\textbf{\S\ref{sec:rw-kgrec} KG recommenders elsewhere}\\
  mature mechanisms exist, but none has been applied to skincare};
\node[box, below=of c] (d) {\textbf{\S\ref{sec:rw-dl} Deep learning}\\
  effective where labelled data exists, cannot cite a regulator, needs interactions};
\node[box, below=of d] (e) {\textbf{\S\ref{sec:rw-hybrid} Hybrid neuro-symbolic}\\
  assumes a domain knowledge graph already exists};
\node[gap,  below=of e] (f) {\S\ref{sec:rw-gap} Gap: no skincare knowledge graph with
  regulator-linked ingredients, modelled provenance and real market availability};
\draw[arr] (a) -- (b); \draw[arr] (b) -- (c);
\draw[arr] (c) -- (d); \draw[arr] (d) -- (e); \draw[arr] (e) -- (f);
\end{tikzpicture}
\caption{The structure of this chapter. Each section identifies a limitation
that the next addresses.}
\label{fig:rw-structure}
\end{figure}

% ---------------------------------------------------------------------
\section{Dataset and resource construction}
\label{sec:rw-datasets}

% TODO Lynne: paste your existing dataset related work from main.tex here,
% then keep the closing paragraph below, which is what links this section
% to the next one.

Resource contributions are an established category with their own venues and
their own review criteria. FoodKG was published as a resource paper at the
International Semantic Web Conference~\cite{haussmann2019foodkg}, and the TOXIN
knowledge graph appeared in \emph{Database}~\cite{toxin2024}. Both are judged on
availability, documentation, validation and a stated maintenance plan rather
than on predictive accuracy.

% TODO Lynne: one paragraph here on existing cosmetic product datasets and
% why none covers a national market. Use the material already in main.tex.

\paragraph{What is missing.} Existing cosmetic product resources record what a
product is. None records whether the product can be obtained in a given market,
what it costs there, or how much the source of each claim is worth. The datasets
that do exist are not linked to any regulatory register, so no statement they
support can be traced to a legal instrument. The next section shows that the
ontologies built over such data inherit all three limitations.

% ---------------------------------------------------------------------
\section{Ontologies and knowledge graphs for skincare and cosmetics}
\label{sec:rw-ontologies}

\subsection{Cross-domain cosmetics ontologies}
\label{ssec:rw-moe}

Moe and Aung~\cite{moe2014} construct two ontologies, one describing facial
skin problems and one describing cosmetics, joined by bridge relations. Their
methodological contribution is an explicit eight-task ontological engineering
procedure, running from a glossary of terms with synonyms and acronyms through
concept taxonomies, binary relation diagrams and a concept dictionary to a
constants table. Both ontologies were built in Prot\'eg\'e.

Recommendation proceeds in three stages. Taxonomic conversational case-based
reasoning narrows an initially vague user query through ranked questions, using
an \texttt{isNextRelatedTo} property to order them. The resulting problem and
the candidate products are then placed in a weighted directed acyclic graph, and
the Ford--Fulkerson maximum flow algorithm assigns each product a flow weight
$W(v_i)=\sum_k f(v_i,k)$, with products ranked in descending order of that
weight. Evaluation reports precision, recall and $F$-measure as a function of a
recommendation threshold $\alpha$.

Two observations matter for this thesis. First, no reasoner is used at any
point, so the ontology serves as a graph structure rather than as a logical
theory. Second, the paper claims greater accuracy than related work without
reporting a dataset size, a baseline or a comparison, so that claim is not
supported by the evidence presented. Ingredients are represented as a free-text
property with an associated numeric value and are not linked to any register,
which means no statement about regulatory status is expressible in the model.

\subsection{OntoCosmetic and formulation support}
\label{ssec:rw-ontocosmetic}

Serna et al.~\cite{serna2021} present OntoCosmetic, a decision support ontology
for the design of emulsion-based cosmetic products, extended by Gabriel et
al.~\cite{gabriel2023} into a cross-platform application, Formultools. The
knowledge base is assembled from four kinds of material: general subproblems
expressed as physicochemical properties to promote or limit, general solution
strategies, ingredient databases typed by function, and heuristics linking the
two. Heuristics are encoded as SWRL rules in Prot\'eg\'e and carry
\texttt{hasHeuristicHighThreshold}, \texttt{hasHeuristicLowThreshold} and
\texttt{hasHeuristicSource} properties, the last of which records the provenance
of the rule itself.

The application was developed using the five-planes user-centred design
method, in co-design with cosmetics experts, and evaluated across iterations
using the AttrakDiff instrument. It supports ingredient screening, multi-criteria
ingredient selection and formulation evaluation, citing the Analytic Hierarchy
Process for the multi-criteria component.

OntoCosmetic is the only ontology reviewed in this section that is publicly
resolvable, at \url{https://purl.org/ontocosmetic}. Inspection of the published
file gives 116 classes, 26 object properties, 20 data properties and 279
individuals, all populated manually. Its classes, among them \texttt{HLB},
\texttt{DropletSize}, \texttt{Rheology} and \texttt{Emollient\_Dosage}, model
formulation chemistry rather than retail selection: the ontology is designed to
help a formulator make a stable cream, not to help a consumer choose one.
Neither paper reports an evaluation against ground truth, and Gabriel et al.
state in their conclusion that testing with experts remained future work.

\subsection{Vision-driven skincare recommendation}
\label{ssec:rw-hansanie}

Hansanie and Silva~\cite{hansanie2024} combine a convolutional neural network
that grades acne severity from a facial photograph with an ontology that
performs the product selection. Their vocabulary was derived from interviews
with dermatologists and health professionals together with a survey of
twenty-one participants, and the class hierarchy was built top-down. The
ontology is deliberately split into three modules, covering skincare concepts,
product information and user profiles, which are then merged; the stated
rationale is that the three evolve independently.

Their property set includes \texttt{suitableFor} from treatment products to
skin types, \texttt{hasKeyIngredient}, \texttt{hasProductRecommendation} and
\texttt{hasRating}. The implementation uses Prot\'eg\'e, the Pellet reasoner,
Owlready2 for querying from Python, and Tkinter for the interface. The severity
grade produced by the network is asserted into the ontology as a fact about the
user, so the perception component and the knowledge component remain separable.

The reported figures require care. The headline result of 87.5\,\% is the
proportion of twenty-four survey participants who expressed satisfaction with
the products shown; the classification accuracy of the network itself is
77.5\,\%. No baseline recommender is compared against, and no ground-truth set
of correct recommendations is constructed.

\subsection{Defined classes and query-driven recommendation}
\label{ssec:rw-abesova}

Abesov\'a et al.~\cite{abesova2023} build a skincare ontology over Sephora
product and review data, comprising 36 classes, 6 object properties and 6 data
properties, following an explicitly middle-out and iterative construction
process in which the scope widened twice as new data sources were incorporated.
The knowledge graph was produced from scraped CSV files using OntoRefine and
stored in GraphDB, with country data and capital-city coordinates obtained by
querying DBpedia through an external SPARQL endpoint. Recommendation is
performed by SPARQL queries rather than by procedural code.

Their principal methodological contribution, and the one adopted in this
thesis, is the use of defined classes. Rather than annotating individual
products with a suitability label, they state an equivalence axiom, so that
membership of \texttt{OilySkinProducts} follows from
\texttt{hasOilyScore\ value\ 5} and is recomputed by the reasoner whenever the
underlying data changes.

The work is a student project and is not peer reviewed, which is stated here
because it is cited for a mechanism rather than for a result. Its three
acknowledged limitations are directly relevant. Defined-class membership
provides no ordering within a class, so ranking requires a separate signal; a
product supported by a single weak review can be recommended ahead of products
with strong support, because the two selection paths are not commensurable; and
the authors observe that class restrictions based on ingredients would permit a
symbolic rather than statistical approach, which they did not implement.

\subsection{Ontologies populated from commercial catalogues}
\label{ssec:rw-bittech}

A recent ontology-based skincare recommendation system~\cite{bittech2025}
follows METHONTOLOGY to produce twelve core classes, among them Product,
Ingredient, Skin Type, Skin Concern, Benefit, Allergen Type and Formulation
Trait, together with more than twenty-five object properties. The knowledge
base was populated by scraping three commercial platforms, yielding over 3{,}800
products and 28{,}000 ingredients, and is queried through Apache Jena Fuseki.

This is the closest published system to the present work in both scale and
intent. It differs in four respects. It targets the Indonesian market rather
than a market in which product availability is itself uncertain; its allergen
representation is a locally defined class rather than a link to a statutory
register; it does not model price, availability or provenance; and its product
count is approximately one third of the dataset used here.

% TODO Lynne: obtain and read the full PDF before submission, then confirm or
% correct the four points above and add whether their ontology is published.

\subsection{Regulatory and clinical knowledge graphs}
\label{ssec:rw-regulatory}

Two lines of work outside recommendation are relevant because they supply the
grounding that skincare ontologies lack.

The TOXIN knowledge graph~\cite{toxin2024} supports animal-free risk assessment
of cosmetics, drawing on the scientific opinions that underpin the annexes of
Regulation (EC) No~1223/2009. Its construction is instructive: curated
spreadsheets are transformed into RDF using R2RML, the ToXic Process Ontology
from the OBO Foundry provides the primary structure, external resources are
incorporated by reference to their IRIs rather than by copying, and each
integration is held in a separate named graph so that the origin of any
statement can be traced. It contains no products.

In the clinical direction, DermO~\cite{dermo2016} provides more than three
thousand expert-authored terms for cutaneous disease, organised under twenty
upper-level categories, aligned to ICD-10 and published through BioPortal in
both OBO and OWL~2 formats. The more recent D3X ontology~\cite{d3x2024} links
dermoscopic patterns to differential diagnoses and is evaluated through a
usability study. Neither addresses cosmetic products, but both offer a
vocabulary for skin concerns with clinical standing that a retail-derived
vocabulary does not possess.

A parallel may also be drawn with ontologies that model regulatory status
directly. Work on flavouring ontologies for halal status determination
\cite{halal2024} represents the ingredient, the certifying authority and the
resulting status as distinct entities, so that a product inherits a status from
its constituents according to a named authority. The structure transfers
directly to ingredients restricted under the European cosmetics regulation.

\subsection{Summary of Section~\ref{sec:rw-ontologies}}

\begin{table}[htbp]
\centering
\caption{Ontologies and knowledge graphs in skincare, cosmetics and cosmetic
regulation. A dash indicates the property is absent or not reported.}
\label{tab:rw-ontologies}
\footnotesize
\setlength{\tabcolsep}{4pt}
\begin{tabular}{lccccccc}
\toprule
 & \rotatebox{60}{Moe 2014} & \rotatebox{60}{OntoCosmetic}
 & \rotatebox{60}{Hansanie 2024} & \rotatebox{60}{Abesov\'a 2023}
 & \rotatebox{60}{Bit-Tech 2025} & \rotatebox{60}{TOXIN 2024}
 & \rotatebox{60}{\textbf{This work}} \\
\midrule
Peer reviewed                & $\circ$ & \checkmark & \checkmark & --         & \checkmark & \checkmark & \checkmark \\
Ontology published           & --      & \checkmark & --         & --         & --         & \checkmark & \checkmark \\
Named methodology            & $\circ$ & --         & --         & \checkmark & \checkmark & --         & \checkmark \\
Products modelled            & --      & 279        & $\circ$    & $\circ$    & 3{,}800    & --         & \textbf{12{,}629} \\
Ingredients linked to registry & --    & --         & --         & --         & --         & \checkmark & \checkmark \\
Declarative population       & --      & --         & --         & \checkmark & --         & \checkmark & \checkmark \\
Reasoner used                & --      & \checkmark & \checkmark & \checkmark & \checkmark & --         & \checkmark \\
SPARQL interface             & --      & --         & $\circ$    & \checkmark & \checkmark & \checkmark & \checkmark \\
External linking             & --      & --         & --         & \checkmark & --         & \checkmark & \checkmark \\
Provenance modelled          & --      & $\circ$    & --         & --         & --         & \checkmark & \checkmark \\
Price or availability        & $\circ$ & $\circ$    & --         & $\circ$    & --         & --         & \checkmark \\
Evaluated on ground truth    & --      & --         & $\circ$    & --         & $\circ$    & n/a        & \checkmark \\
\bottomrule
\end{tabular}
\begin{tablenotes}\footnotesize
\item \checkmark\ present; $\circ$ partial or weak; -- absent or not reported.
\end{tablenotes}
\end{table}

\paragraph{What is missing.} No system in Table~\ref{tab:rw-ontologies}
combines an ingredient layer with statutory grounding, an explicit model of
where each claim came from, and a representation of whether the recommended
product can actually be obtained. Ingredients are represented as free text in
every skincare ontology reviewed; the single work that links ingredients to the
European regulatory framework contains no products at all. The mechanisms
needed to close this gap have, however, been developed in other domains, and
these are reviewed next.

% ---------------------------------------------------------------------
\section{Ontology and knowledge graph based recommender systems}
\label{sec:rw-kgrec}

\subsection{Ontological user profiling}

Middleton et al.~\cite{middleton2004} established ontological user profiling in
their Quickstep and Foxtrot systems for recommending research papers, reporting
three findings that remain the standard justification for the approach:
ontological inference improves user profiling, because interest in a specific
topic implies interest in its parent; external ontological knowledge can
bootstrap a recommender for users with no history; and visualising a profile so
that users may correct it improves accuracy. The second finding is the direct
answer to the cold-start problem that affects any purely collaborative method,
and it applies with particular force to a market in which no interaction logs
exist.

\subsection{Constraint satisfaction over a domain knowledge graph}

FoodKG~\cite{haussmann2019foodkg} integrates recipes, nutritional data, food
taxonomies and links to existing ontologies into a single graph, over which a
SPARQL service identifies recipes that can be prepared from available
ingredients while respecting hard constraints such as allergies. The structural
correspondence with the present work is close: recipes correspond to products,
nutritional data supplied by an authority corresponds to the CosIng register,
and allergen constraints correspond to the twenty-six fragrance allergens whose
declaration is required under the European regulation. Subsequent work frames
the task explicitly as constrained question answering over a knowledge
graph~\cite{chen2021constrained}, which is a more accurate description of the
skincare selection problem than ranking is.

\subsection{Linked open data as a source of item features}

Di Noia, Ostuni and colleagues~\cite{dinoia2012lod,ostuni2013topn,dinoia2015wims}
compute item similarity using a semantic vector space model whose dimensions are
links in the linked open data cloud rather than terms in a document. The
approach permits content-based recommendation in the absence of ratings, which
is the situation of this thesis, where the rating field is populated for only
half of the catalogue and no user history exists at all. The authors identify
entity resolution against external datasets as their principal open problem, an
observation consistent with the matching error rates measured during the
construction of the dataset used here.

\subsection{User needs as first-class entities}

AliCoCo~\cite{alicoco2020} argues that conventional product ontologies leave a
semantic gap because they describe what a product is rather than what a user
needs, and formalises e-commerce concepts, that is user needs, as first-class
entities in their own right. The argument transfers directly. Skin concerns are
presently modelled throughout the literature as attributes of products, whereas
a need entity can be satisfied by a combination of products, can carry a budget
and a season, and can be constrained by local availability.

\subsection{Taxonomies of knowledge graph recommendation}

Guo et al.~\cite{guo2022survey} classify knowledge-graph-based recommender
systems as embedding-based, path-based or unified. All three families presuppose
a user--item interaction matrix and treat the knowledge graph as side
information mitigating sparsity. The present work is situated upstream of this
taxonomy: it constructs, for a domain and a market in which none exists, the
knowledge graph that such methods presuppose.

\paragraph{What is missing.} The mechanisms surveyed in this section, defined
classes, hard constraint satisfaction, linked-data-derived similarity and needs
as entities, are mature and well evidenced, but none has been applied to
skincare, and none of the systems reviewed operates in a market where
availability cannot be assumed.

% ---------------------------------------------------------------------
\section{Deep learning approaches}
\label{sec:rw-dl}

Lee et al.~\cite{lee2024} estimate the efficacy of a cosmetic product from its
ingredient list using a deep neural network, combining the result with automated
facial skin analysis to produce personalised recommendations. The point of
departure is the same as that of this thesis, namely the ingredient list, and
the divergence lies entirely in how ingredients are connected to outcomes: by a
function learned from data in their case, and by the European register together
with stated rules in this one.

The broader literature on neural recommendation is well
surveyed~\cite{zhang2019survey}, with neural collaborative
filtering~\cite{he2017ncf} providing the standard baseline. Every method in this
family learns from user behaviour. The dataset constructed for this thesis
contains none, because it is a product resource rather than an interaction log,
and this determines which half of the literature the work belongs to.

It should be stated directly that where sufficient labelled data exists, a
learned model is likely to outperform a rule-based one on benchmark accuracy,
because it can capture interactions that are not explicitly encoded. The
argument advanced here is not one of accuracy. It is that in a domain carrying
physical risk and governed by a legal framework, a recommendation that cannot
state its reason cannot be audited by a regulator, cannot be used by a
pharmacist, and cannot be corrected when it is wrong.

\paragraph{What is missing.} Vision and learned models can establish the state
of a user's skin and can estimate what a formulation is likely to do. Neither
can establish that the resulting product is stocked in Beirut, what it costs
there this month, or which annex of the cosmetics regulation restricts one of
its preservatives.

% ---------------------------------------------------------------------
\section{Hybrid neuro-symbolic approaches}
\label{sec:rw-hybrid}

A substantial body of work combines knowledge graphs with neural
recommendation. CKE~\cite{zhang2016cke} learns item representations jointly from
a knowledge graph and from ratings; RippleNet~\cite{wang2018ripplenet} propagates
user preference outward through the graph; KGAT~\cite{wang2019kgat} applies graph
attention over a collaborative knowledge graph to capture high-order relations;
and KPRN~\cite{wang2019kprn} encodes the path between user and item with a
recurrent network and weights paths by their strength, so that the path itself
constitutes an explanation. Subsequent work optimises explanation paths for
quality rather than relevance alone~\cite{balloccu2022}.

Each of these methods requires an interaction matrix and treats the knowledge
graph as auxiliary. The position of this thesis is the converse: the graph is
the contribution and the interactions do not yet exist. These methods are
therefore what the resource constructed here enables once deployment produces
interaction data, and they are discussed again in the future work of
Chapter~\ref{ch:conclusion}.

% TODO Lynne: check that \ref{ch:conclusion} matches your actual label.

One bridge is applicable immediately. Ontology embedding methods such as
RDF2Vec~\cite{ristoski2016rdf2vec} and OWL2Vec*~\cite{chen2021owl2vec} require an
ontology rather than an interaction log. OWL2Vec* in particular encodes graph
structure, lexical content and logical constructors, having first materialised
the ontology with a reasoner, and was evaluated on class membership and class
subsumption prediction. Both tasks correspond to real gaps in the present
dataset, in which size is recorded for 27\,\% of products and rating for 51\,\%.

Finally, the use of large language models in ontology engineering is developing
rapidly, with a dedicated shared task~\cite{llms4ol2024} covering term typing,
taxonomy discovery and non-taxonomic relation extraction, and a recent
systematic review of the area~\cite{li2026llmoe}. The free-text benefit and
concern fields of the dataset used here are candidates for exactly these tasks,
subject to validation of every proposed term against an authoritative source.

% ---------------------------------------------------------------------
\section{The gap addressed by this thesis}
\label{sec:rw-gap}

Five limitations recur across the reviewed literature.

\begin{enumerate}
  \item \textbf{No ingredient layer with statutory standing.} Every skincare
    ontology reviewed represents ingredients as free text or as a small,
    manually curated class list. The one knowledge graph that links cosmetic
    ingredients to the European regulatory framework contains no products.
  \item \textbf{No model of provenance.} No reviewed skincare ontology records
    who asserted a claim, or how much weight that source should carry. Every
    claim is presented as equally reliable.
  \item \textbf{No model of availability.} Every reviewed system assumes that a
    recommended product can be purchased. In the Lebanese market that assumption
    fails, and it fails differently for imported and for locally manufactured
    products.
  \item \textbf{Weak evaluation.} Two of the reviewed systems report user
    satisfaction on samples of fewer than thirty participants, one presents a
    single case study, and one states that expert evaluation was never carried
    out.
  \item \textbf{Poor availability of artefacts.} One of six reviewed ontologies
    is publicly resolvable, which prevents the field from building on its own
    results.
\end{enumerate}

This thesis addresses limitations 1, 2, 3 and 5 directly through the design of
the ontology and the publication of the resulting artefacts, and addresses
limitation 4 through a competency-question-driven evaluation reported in
Chapter~\ref{ch:evaluation}.

% TODO Lynne: check that \ref{ch:evaluation} matches your actual label.

%  Requires in the preamble:
%     \usepackage{booktabs}
%     \usepackage{graphicx}
%     \usepackage{tikz}
%     \usetikzlibrary{positioning,arrows.meta,shapes.geometric}
%     \usepackage[hidelinks]{hyperref}
%     \usepackage{multirow}
%     \usepackage{rotating}      % for the wide comparison table
%     \usepackage{threeparttable}
%
%  Bibliography: related_work.bib
% =====================================================================

\chapter{Related Work}
\label{ch:related}

% ---------------------------------------------------------------------
\section{How this chapter is organised}
\label{sec:rw-organisation}

This chapter narrows in four steps. Section~\ref{sec:rw-datasets} reviews how
product datasets and semantic resources are constructed and what the community
expects of them, which positions the first contribution of this thesis.
Section~\ref{sec:rw-ontologies} reviews ontologies and knowledge graphs built
for skincare, cosmetics and cosmetic regulation, which are the direct
competitors. Section~\ref{sec:rw-kgrec} reviews ontology and knowledge graph
based recommender systems in other domains, from which the mechanisms of this
work are drawn. Section~\ref{sec:rw-dl} reviews deep learning approaches,
including those applied to skin analysis, and states plainly what they do
better. Section~\ref{sec:rw-hybrid} reviews hybrid neuro-symbolic work and
positions it as future rather than present work. Section~\ref{sec:rw-gap}
consolidates the gap.

Each section closes by naming what is missing, and the missing thing is what
the following section addresses. The final gap is the thesis.

\begin{figure}[htbp]
\centering
\begin{tikzpicture}[
  node distance=7mm,
  box/.style={draw, rounded corners, align=center, text width=0.82\linewidth,
              inner sep=5pt, font=\small},
  gap/.style={draw, thick, rounded corners, align=center,
              text width=0.82\linewidth, inner sep=5pt, font=\small\bfseries},
  arr/.style={-{Latex[length=2mm]}, thick}
]
\node[box] (a) {\textbf{\S\ref{sec:rw-datasets} Resource construction}\\
  provenance, validation and availability are expected of a modern resource};
\node[box, below=of a] (b) {\textbf{\S\ref{sec:rw-ontologies} Skincare and cosmetics ontologies}\\
  small, mostly unpublished, ingredients as free text, no regulatory grounding};
\node[box, below=of b] (c) {\textbf{\S\ref{sec:rw-kgrec} KG recommenders elsewhere}\\
  mature mechanisms exist, but none has been applied to skincare};
\node[box, below=of c] (d) {\textbf{\S\ref{sec:rw-dl} Deep learning}\\
  effective where labelled data exists, cannot cite a regulator, needs interactions};
\node[box, below=of d] (e) {\textbf{\S\ref{sec:rw-hybrid} Hybrid neuro-symbolic}\\
  assumes a domain knowledge graph already exists};
\node[gap,  below=of e] (f) {\S\ref{sec:rw-gap} Gap: no skincare knowledge graph with
  regulator-linked ingredients, modelled provenance and real market availability};
\draw[arr] (a) -- (b); \draw[arr] (b) -- (c);
\draw[arr] (c) -- (d); \draw[arr] (d) -- (e); \draw[arr] (e) -- (f);
\end{tikzpicture}
\caption{The structure of this chapter. Each section identifies a limitation
that the next addresses.}
\label{fig:rw-structure}
\end{figure}

% ---------------------------------------------------------------------
\section{Dataset and resource construction}
\label{sec:rw-datasets}

% TODO Lynne: paste your existing dataset related work from main.tex here,
% then keep the closing paragraph below, which is what links this section
% to the next one.

Resource contributions are an established category with their own venues and
their own review criteria. FoodKG was published as a resource paper at the
International Semantic Web Conference~\cite{haussmann2019foodkg}, and the TOXIN
knowledge graph appeared in \emph{Database}~\cite{toxin2024}. Both are judged on
availability, documentation, validation and a stated maintenance plan rather
than on predictive accuracy.

% TODO Lynne: one paragraph here on existing cosmetic product datasets and
% why none covers a national market. Use the material already in main.tex.

\paragraph{What is missing.} Existing cosmetic product resources record what a
product is. None records whether the product can be obtained in a given market,
what it costs there, or how much the source of each claim is worth. The datasets
that do exist are not linked to any regulatory register, so no statement they
support can be traced to a legal instrument. The next section shows that the
ontologies built over such data inherit all three limitations.

% ---------------------------------------------------------------------
\section{Ontologies and knowledge graphs for skincare and cosmetics}
\label{sec:rw-ontologies}

\subsection{Cross-domain cosmetics ontologies}
\label{ssec:rw-moe}

Moe and Aung~\cite{moe2014} construct two ontologies, one describing facial
skin problems and one describing cosmetics, joined by bridge relations. Their
methodological contribution is an explicit eight-task ontological engineering
procedure, running from a glossary of terms with synonyms and acronyms through
concept taxonomies, binary relation diagrams and a concept dictionary to a
constants table. Both ontologies were built in Prot\'eg\'e.

Recommendation proceeds in three stages. Taxonomic conversational case-based
reasoning narrows an initially vague user query through ranked questions, using
an \texttt{isNextRelatedTo} property to order them. The resulting problem and
the candidate products are then placed in a weighted directed acyclic graph, and
the Ford--Fulkerson maximum flow algorithm assigns each product a flow weight
$W(v_i)=\sum_k f(v_i,k)$, with products ranked in descending order of that
weight. Evaluation reports precision, recall and $F$-measure as a function of a
recommendation threshold $\alpha$.

Two observations matter for this thesis. First, no reasoner is used at any
point, so the ontology serves as a graph structure rather than as a logical
theory. Second, the paper claims greater accuracy than related work without
reporting a dataset size, a baseline or a comparison, so that claim is not
supported by the evidence presented. Ingredients are represented as a free-text
property with an associated numeric value and are not linked to any register,
which means no statement about regulatory status is expressible in the model.

\subsection{OntoCosmetic and formulation support}
\label{ssec:rw-ontocosmetic}

Serna et al.~\cite{serna2021} present OntoCosmetic, a decision support ontology
for the design of emulsion-based cosmetic products, extended by Gabriel et
al.~\cite{gabriel2023} into a cross-platform application, Formultools. The
knowledge base is assembled from four kinds of material: general subproblems
expressed as physicochemical properties to promote or limit, general solution
strategies, ingredient databases typed by function, and heuristics linking the
two. Heuristics are encoded as SWRL rules in Prot\'eg\'e and carry
\texttt{hasHeuristicHighThreshold}, \texttt{hasHeuristicLowThreshold} and
\texttt{hasHeuristicSource} properties, the last of which records the provenance
of the rule itself.

The application was developed using the five-planes user-centred design
method, in co-design with cosmetics experts, and evaluated across iterations
using the AttrakDiff instrument. It supports ingredient screening, multi-criteria
ingredient selection and formulation evaluation, citing the Analytic Hierarchy
Process for the multi-criteria component.

OntoCosmetic is the only ontology reviewed in this section that is publicly
resolvable, at \url{https://purl.org/ontocosmetic}. Inspection of the published
file gives 116 classes, 26 object properties, 20 data properties and 279
individuals, all populated manually. Its classes, among them \texttt{HLB},
\texttt{DropletSize}, \texttt{Rheology} and \texttt{Emollient\_Dosage}, model
formulation chemistry rather than retail selection: the ontology is designed to
help a formulator make a stable cream, not to help a consumer choose one.
Neither paper reports an evaluation against ground truth, and Gabriel et al.
state in their conclusion that testing with experts remained future work.

\subsection{Vision-driven skincare recommendation}
\label{ssec:rw-hansanie}

Hansanie and Silva~\cite{hansanie2024} combine a convolutional neural network
that grades acne severity from a facial photograph with an ontology that
performs the product selection. Their vocabulary was derived from interviews
with dermatologists and health professionals together with a survey of
twenty-one participants, and the class hierarchy was built top-down. The
ontology is deliberately split into three modules, covering skincare concepts,
product information and user profiles, which are then merged; the stated
rationale is that the three evolve independently.

Their property set includes \texttt{suitableFor} from treatment products to
skin types, \texttt{hasKeyIngredient}, \texttt{hasProductRecommendation} and
\texttt{hasRating}. The implementation uses Prot\'eg\'e, the Pellet reasoner,
Owlready2 for querying from Python, and Tkinter for the interface. The severity
grade produced by the network is asserted into the ontology as a fact about the
user, so the perception component and the knowledge component remain separable.

The reported figures require care. The headline result of 87.5\,\% is the
proportion of twenty-four survey participants who expressed satisfaction with
the products shown; the classification accuracy of the network itself is
77.5\,\%. No baseline recommender is compared against, and no ground-truth set
of correct recommendations is constructed.

\subsection{Defined classes and query-driven recommendation}
\label{ssec:rw-abesova}

Abesov\'a et al.~\cite{abesova2023} build a skincare ontology over Sephora
product and review data, comprising 36 classes, 6 object properties and 6 data
properties, following an explicitly middle-out and iterative construction
process in which the scope widened twice as new data sources were incorporated.
The knowledge graph was produced from scraped CSV files using OntoRefine and
stored in GraphDB, with country data and capital-city coordinates obtained by
querying DBpedia through an external SPARQL endpoint. Recommendation is
performed by SPARQL queries rather than by procedural code.

Their principal methodological contribution, and the one adopted in this
thesis, is the use of defined classes. Rather than annotating individual
products with a suitability label, they state an equivalence axiom, so that
membership of \texttt{OilySkinProducts} follows from
\texttt{hasOilyScore\ value\ 5} and is recomputed by the reasoner whenever the
underlying data changes.

The work is a student project and is not peer reviewed, which is stated here
because it is cited for a mechanism rather than for a result. Its three
acknowledged limitations are directly relevant. Defined-class membership
provides no ordering within a class, so ranking requires a separate signal; a
product supported by a single weak review can be recommended ahead of products
with strong support, because the two selection paths are not commensurable; and
the authors observe that class restrictions based on ingredients would permit a
symbolic rather than statistical approach, which they did not implement.

\subsection{Ontologies populated from commercial catalogues}
\label{ssec:rw-bittech}

A recent ontology-based skincare recommendation system~\cite{bittech2025}
follows METHONTOLOGY to produce twelve core classes, among them Product,
Ingredient, Skin Type, Skin Concern, Benefit, Allergen Type and Formulation
Trait, together with more than twenty-five object properties. The knowledge
base was populated by scraping three commercial platforms, yielding over 3{,}800
products and 28{,}000 ingredients, and is queried through Apache Jena Fuseki.

This is the closest published system to the present work in both scale and
intent. It differs in four respects. It targets the Indonesian market rather
than a market in which product availability is itself uncertain; its allergen
representation is a locally defined class rather than a link to a statutory
register; it does not model price, availability or provenance; and its product
count is approximately one third of the dataset used here.

% TODO Lynne: obtain and read the full PDF before submission, then confirm or
% correct the four points above and add whether their ontology is published.

\subsection{Regulatory and clinical knowledge graphs}
\label{ssec:rw-regulatory}

Two lines of work outside recommendation are relevant because they supply the
grounding that skincare ontologies lack.

The TOXIN knowledge graph~\cite{toxin2024} supports animal-free risk assessment
of cosmetics, drawing on the scientific opinions that underpin the annexes of
Regulation (EC) No~1223/2009. Its construction is instructive: curated
spreadsheets are transformed into RDF using R2RML, the ToXic Process Ontology
from the OBO Foundry provides the primary structure, external resources are
incorporated by reference to their IRIs rather than by copying, and each
integration is held in a separate named graph so that the origin of any
statement can be traced. It contains no products.

In the clinical direction, DermO~\cite{dermo2016} provides more than three
thousand expert-authored terms for cutaneous disease, organised under twenty
upper-level categories, aligned to ICD-10 and published through BioPortal in
both OBO and OWL~2 formats. The more recent D3X ontology~\cite{d3x2024} links
dermoscopic patterns to differential diagnoses and is evaluated through a
usability study. Neither addresses cosmetic products, but both offer a
vocabulary for skin concerns with clinical standing that a retail-derived
vocabulary does not possess.

A parallel may also be drawn with ontologies that model regulatory status
directly. Work on flavouring ontologies for halal status determination
\cite{halal2024} represents the ingredient, the certifying authority and the
resulting status as distinct entities, so that a product inherits a status from
its constituents according to a named authority. The structure transfers
directly to ingredients restricted under the European cosmetics regulation.

\subsection{Summary of Section~\ref{sec:rw-ontologies}}

\begin{table}[htbp]
\centering
\caption{Ontologies and knowledge graphs in skincare, cosmetics and cosmetic
regulation. A dash indicates the property is absent or not reported.}
\label{tab:rw-ontologies}
\footnotesize
\setlength{\tabcolsep}{4pt}
\begin{tabular}{lccccccc}
\toprule
 & \rotatebox{60}{Moe 2014} & \rotatebox{60}{OntoCosmetic}
 & \rotatebox{60}{Hansanie 2024} & \rotatebox{60}{Abesov\'a 2023}
 & \rotatebox{60}{Bit-Tech 2025} & \rotatebox{60}{TOXIN 2024}
 & \rotatebox{60}{\textbf{This work}} \\
\midrule
Peer reviewed                & $\circ$ & \checkmark & \checkmark & --         & \checkmark & \checkmark & \checkmark \\
Ontology published           & --      & \checkmark & --         & --         & --         & \checkmark & \checkmark \\
Named methodology            & $\circ$ & --         & --         & \checkmark & \checkmark & --         & \checkmark \\
Products modelled            & --      & 279        & $\circ$    & $\circ$    & 3{,}800    & --         & \textbf{12{,}629} \\
Ingredients linked to registry & --    & --         & --         & --         & --         & \checkmark & \checkmark \\
Declarative population       & --      & --         & --         & \checkmark & --         & \checkmark & \checkmark \\
Reasoner used                & --      & \checkmark & \checkmark & \checkmark & \checkmark & --         & \checkmark \\
SPARQL interface             & --      & --         & $\circ$    & \checkmark & \checkmark & \checkmark & \checkmark \\
External linking             & --      & --         & --         & \checkmark & --         & \checkmark & \checkmark \\
Provenance modelled          & --      & $\circ$    & --         & --         & --         & \checkmark & \checkmark \\
Price or availability        & $\circ$ & $\circ$    & --         & $\circ$    & --         & --         & \checkmark \\
Evaluated on ground truth    & --      & --         & $\circ$    & --         & $\circ$    & n/a        & \checkmark \\
\bottomrule
\end{tabular}
\begin{tablenotes}\footnotesize
\item \checkmark\ present; $\circ$ partial or weak; -- absent or not reported.
\end{tablenotes}
\end{table}

\paragraph{What is missing.} No system in Table~\ref{tab:rw-ontologies}
combines an ingredient layer with statutory grounding, an explicit model of
where each claim came from, and a representation of whether the recommended
product can actually be obtained. Ingredients are represented as free text in
every skincare ontology reviewed; the single work that links ingredients to the
European regulatory framework contains no products at all. The mechanisms
needed to close this gap have, however, been developed in other domains, and
these are reviewed next.

% ---------------------------------------------------------------------
\section{Ontology and knowledge graph based recommender systems}
\label{sec:rw-kgrec}

\subsection{Ontological user profiling}

Middleton et al.~\cite{middleton2004} established ontological user profiling in
their Quickstep and Foxtrot systems for recommending research papers, reporting
three findings that remain the standard justification for the approach:
ontological inference improves user profiling, because interest in a specific
topic implies interest in its parent; external ontological knowledge can
bootstrap a recommender for users with no history; and visualising a profile so
that users may correct it improves accuracy. The second finding is the direct
answer to the cold-start problem that affects any purely collaborative method,
and it applies with particular force to a market in which no interaction logs
exist.

\subsection{Constraint satisfaction over a domain knowledge graph}

FoodKG~\cite{haussmann2019foodkg} integrates recipes, nutritional data, food
taxonomies and links to existing ontologies into a single graph, over which a
SPARQL service identifies recipes that can be prepared from available
ingredients while respecting hard constraints such as allergies. The structural
correspondence with the present work is close: recipes correspond to products,
nutritional data supplied by an authority corresponds to the CosIng register,
and allergen constraints correspond to the twenty-six fragrance allergens whose
declaration is required under the European regulation. Subsequent work frames
the task explicitly as constrained question answering over a knowledge
graph~\cite{chen2021constrained}, which is a more accurate description of the
skincare selection problem than ranking is.

\subsection{Linked open data as a source of item features}

Di Noia, Ostuni and colleagues~\cite{dinoia2012lod,ostuni2013topn,dinoia2015wims}
compute item similarity using a semantic vector space model whose dimensions are
links in the linked open data cloud rather than terms in a document. The
approach permits content-based recommendation in the absence of ratings, which
is the situation of this thesis, where the rating field is populated for only
half of the catalogue and no user history exists at all. The authors identify
entity resolution against external datasets as their principal open problem, an
observation consistent with the matching error rates measured during the
construction of the dataset used here.

\subsection{User needs as first-class entities}

AliCoCo~\cite{alicoco2020} argues that conventional product ontologies leave a
semantic gap because they describe what a product is rather than what a user
needs, and formalises e-commerce concepts, that is user needs, as first-class
entities in their own right. The argument transfers directly. Skin concerns are
presently modelled throughout the literature as attributes of products, whereas
a need entity can be satisfied by a combination of products, can carry a budget
and a season, and can be constrained by local availability.

\subsection{Taxonomies of knowledge graph recommendation}

Guo et al.~\cite{guo2022survey} classify knowledge-graph-based recommender
systems as embedding-based, path-based or unified. All three families presuppose
a user--item interaction matrix and treat the knowledge graph as side
information mitigating sparsity. The present work is situated upstream of this
taxonomy: it constructs, for a domain and a market in which none exists, the
knowledge graph that such methods presuppose.

\paragraph{What is missing.} The mechanisms surveyed in this section, defined
classes, hard constraint satisfaction, linked-data-derived similarity and needs
as entities, are mature and well evidenced, but none has been applied to
skincare, and none of the systems reviewed operates in a market where
availability cannot be assumed.

% ---------------------------------------------------------------------
\section{Deep learning approaches}
\label{sec:rw-dl}

Lee et al.~\cite{lee2024} estimate the efficacy of a cosmetic product from its
ingredient list using a deep neural network, combining the result with automated
facial skin analysis to produce personalised recommendations. The point of
departure is the same as that of this thesis, namely the ingredient list, and
the divergence lies entirely in how ingredients are connected to outcomes: by a
function learned from data in their case, and by the European register together
with stated rules in this one.

The broader literature on neural recommendation is well
surveyed~\cite{zhang2019survey}, with neural collaborative
filtering~\cite{he2017ncf} providing the standard baseline. Every method in this
family learns from user behaviour. The dataset constructed for this thesis
contains none, because it is a product resource rather than an interaction log,
and this determines which half of the literature the work belongs to.

It should be stated directly that where sufficient labelled data exists, a
learned model is likely to outperform a rule-based one on benchmark accuracy,
because it can capture interactions that are not explicitly encoded. The
argument advanced here is not one of accuracy. It is that in a domain carrying
physical risk and governed by a legal framework, a recommendation that cannot
state its reason cannot be audited by a regulator, cannot be used by a
pharmacist, and cannot be corrected when it is wrong.

\paragraph{What is missing.} Vision and learned models can establish the state
of a user's skin and can estimate what a formulation is likely to do. Neither
can establish that the resulting product is stocked in Beirut, what it costs
there this month, or which annex of the cosmetics regulation restricts one of
its preservatives.

% ---------------------------------------------------------------------
\section{Hybrid neuro-symbolic approaches}
\label{sec:rw-hybrid}

A substantial body of work combines knowledge graphs with neural
recommendation. CKE~\cite{zhang2016cke} learns item representations jointly from
a knowledge graph and from ratings; RippleNet~\cite{wang2018ripplenet} propagates
user preference outward through the graph; KGAT~\cite{wang2019kgat} applies graph
attention over a collaborative knowledge graph to capture high-order relations;
and KPRN~\cite{wang2019kprn} encodes the path between user and item with a
recurrent network and weights paths by their strength, so that the path itself
constitutes an explanation. Subsequent work optimises explanation paths for
quality rather than relevance alone~\cite{balloccu2022}.

Each of these methods requires an interaction matrix and treats the knowledge
graph as auxiliary. The position of this thesis is the converse: the graph is
the contribution and the interactions do not yet exist. These methods are
therefore what the resource constructed here enables once deployment produces
interaction data, and they are discussed again in the future work of
Chapter~\ref{ch:conclusion}.

% TODO Lynne: check that \ref{ch:conclusion} matches your actual label.

One bridge is applicable immediately. Ontology embedding methods such as
RDF2Vec~\cite{ristoski2016rdf2vec} and OWL2Vec*~\cite{chen2021owl2vec} require an
ontology rather than an interaction log. OWL2Vec* in particular encodes graph
structure, lexical content and logical constructors, having first materialised
the ontology with a reasoner, and was evaluated on class membership and class
subsumption prediction. Both tasks correspond to real gaps in the present
dataset, in which size is recorded for 27\,\% of products and rating for 51\,\%.

Finally, the use of large language models in ontology engineering is developing
rapidly, with a dedicated shared task~\cite{llms4ol2024} covering term typing,
taxonomy discovery and non-taxonomic relation extraction, and a recent
systematic review of the area~\cite{li2026llmoe}. The free-text benefit and
concern fields of the dataset used here are candidates for exactly these tasks,
subject to validation of every proposed term against an authoritative source.

% ---------------------------------------------------------------------
\section{The gap addressed by this thesis}
\label{sec:rw-gap}

Five limitations recur across the reviewed literature.

\begin{enumerate}
  \item \textbf{No ingredient layer with statutory standing.} Every skincare
    ontology reviewed represents ingredients as free text or as a small,
    manually curated class list. The one knowledge graph that links cosmetic
    ingredients to the European regulatory framework contains no products.
  \item \textbf{No model of provenance.} No reviewed skincare ontology records
    who asserted a claim, or how much weight that source should carry. Every
    claim is presented as equally reliable.
  \item \textbf{No model of availability.} Every reviewed system assumes that a
    recommended product can be purchased. In the Lebanese market that assumption
    fails, and it fails differently for imported and for locally manufactured
    products.
  \item \textbf{Weak evaluation.} Two of the reviewed systems report user
    satisfaction on samples of fewer than thirty participants, one presents a
    single case study, and one states that expert evaluation was never carried
    out.
  \item \textbf{Poor availability of artefacts.} One of six reviewed ontologies
    is publicly resolvable, which prevents the field from building on its own
    results.
\end{enumerate}

This thesis addresses limitations 1, 2, 3 and 5 directly through the design of
the ontology and the publication of the resulting artefacts, and addresses
limitation 4 through a competency-question-driven evaluation reported in
Chapter~\ref{ch:evaluation}.

% TODO Lynne: check that \ref{ch:evaluation} matches your actual label.
